\worldchapter{Pantheon}{pantheon}
\index{Pantheon}
\label{ch:pantheon}

\begin{multicols}{2}

\wbif{tirakan}{
The deity world of Tirakan is diverse and difficult for amateurs to keep track of. All cultures of the world have their own deities, which are more or less present. In general, gods on Tirakan are very approachable, many can be invoked directly. The peoples of Tirakan pray for certain weather, for personal luck, for success in battle, or for fellow humans.

These rules reflect the closeness of the peoples to the world of the gods.
}{
This extension brings the work of gods into your campaign. Characters are able to invoke Divine action and have a \textbf{attitude} and \textbf{grace} with their deity. There are various forms of invocation which can be performed by a believer.

The rulebook deliberately refrains from using earthly gods or beliefs here, but there are no limits to the imagination. For a cultist, for example, a being from the Cthulhu mythos can also be a deity.
}

\section{Level of faith}
\index{Level of faith}

\wbif{tirakan}{
Similar to magic, Tirakan's faith evolves over the centuries. While the churches pray for a long time in silent waiting for the return of the gods' work, the influence of the gods develops into a very strong, direct influence by the end of the age. This is represented by the \textbf{faith level}, which behaves similarly to the \textbf{magic level} and changes over the centuries.
\begin{itemize}
    \item 1st century: Faith level 1
    \item 2nd century: Faith level 1
    \item 3rd century: Faith level 1
    \item 4th century: Faith level 1
    \item 5th century: Faith level 1
    \item 6th century: Faith level 2
    \item 7th century: Faith level 3
    \item 8th century: Faith level 4
    \item 9th century: Faith level 5
    \item 10th century: Faith level 6
\end{itemize}
}{
The power of divine activity depends on the \textit{level of faith}. This is a global value that illustrates the strength of divine activity. In general, it is assumed that the world has a faith level of \textbf{3}.

However, particular places can change the level of faith. For example, invocations may be stronger in a large cathedral. Areas may perhaps be subject to a curse, or otherwise have a lower faith level. The faith level, if it differs from 3, is set by the game master.
}

\section{Grace}
\index{Grace}

As a value, grace represents the relationship between services of the priest and favors of the god. The value is 0 at the beginning and can become negative or positive.

The cost of the favors is subtracted from the grace. Grace points can be gained by the priest through godly actions in the game. It depends very much on the type of deity, with which the priest can rise in the deity's favor.

\section{Relics}
\index{Relics}

Relics have a special role in the churches. They strengthen the bond with the god and help the believer to continue on his path.

Common relics are objects from the possession of saints, but also bones of them. But even a simple object related to the deity can be a low level relic, such as a special stone for a diety of stone. The character can get to a relic in many different ways, but it always requires a consecration.

Relics always have a level, which can range from 1 to 6. A level 1 relic can be an object that a saint once touched, for example. A level 6 relic can be a holy weapon or the bones of a saint.

\section{The forms of invocation}
\index{forms of invocation}

There are four forms of invocation to a god. Each of them is performed differently. Each has a different effort and requests a different favor from the deity.

Common to all forms of invocation is the influence of the environment, the priest's condition, as well as faith level of the world. Thus, the following modifications are added to the \textbf{minimum roll} of each invocation (there are invocations that require multiple rolls).
\begin{itemize}
    \item Grace of the priest: \textbf{-(grace/2)}
    \item The intention of the character does not correspond to the virtues of the deity: \textbf{+10}
    \item Ceremonial design (candles, clean cloths, etc.) not present: \textbf{+5}
    \item The attitude of the priest is contrary to the deity: \textbf{+15}
    \item The request is not the first request of the day: \textbf{+2}
    \item Sacrifice is offered: \textbf{-3}
    \item The priest uses incense: \textbf{-2}
    \item The invocation is done in \wbif{tirakan}{Doldag}{Latin}: \textbf{-2}
    \item The invocation is chanted (additional performance check): \textbf{-5}
    \item The prevailing level of faith: \textbf{-faith level}
    \item Additional priests at the invocation: \textbf{-Number}
    \item Relic present: \textbf{-Level}
\end{itemize}

\subsection{Shock prayer}
\index{Shock prayer}

The least form of request is the Shock Prayer. In a short, pleading invocation of 3 seconds, the priest can gain a bonus to one of his attributes or skills. The bonus is equal to \textbf{faith level} points and lasts for \textbf{faith level} minutes.

A Shock Prayer requires a single \textbf{Charm} roll.

The Shock Prayer costs the priest 2 grace points.

\subsection{Blessing}
\index{Blessing}

A blessing is able to break a divine curse (the work of a dark god, as indicated by the work in each case), but can also be transferred to an object to create a blessed weapon, holy water, or the like. To cast the blessing takes 5 minutes, and it lasts indefinitely.

A Blessing requires a Willpower and a Charm check.

The blessing costs the priest 5 grace points.

\subsection{Lesser request}
\index{Lesser request}

The Lesser Request invokes direct divine action. In it, the abilities of the character's deity and all of its servants that are classified as ``minor'' can be requested. The prayer for the low petition takes about 15 minutes. It can also be done as part of a ceremonial service.

A charm roll is required for the lesser request.

The grace cost of the favor depends on the request and ranges from 2 to 12 points.

\subsection{Invocation}
\index{Invocation}

The invocation requests a deity's work that is classified as ``higher''. Again, both the deity of the character and its servants may be invoked. The invocation requires a larger ceremony and lasts at least 30 minutes. It can also be done as part of a ceremonial service.

The invocation requires 2 charm rolls and a willpower roll.

The grace cost of the invocation depends on the request and ranges from 10 to 25 points.

\begin{quote}
A word about the gods' work. The work of the gods is sometimes described with concrete rules. However, most descriptions remain rather vague. This is to reflect the fact that the nature and workds of the gods are their own business. GMs and players should be open to spontaneous developments when a god or demon intervenes in world events.
\end{quote}

\subsection{Consecration}
\index{Consecration}

With the consecration, an item such as a weapon is given to a god. The divine power ensures that the item is improved (stats plus about 30-50\%), however there is also a chance that the item will be ensouled after the consecration and have some life of its own.

A consecration is a two-hour ceremony during which the deity is invoked three times by means of a charm roll. In addition, a test of strength is required as the item is held for the entire period. Finally, a 50\% chance of ensoulment is thrown.

The consecration costs the priest 7 grace points.

\subsection{Silent prayer}
\index{Silent prayer}

Once per day, the priest may spend one hour in silent devotion to his deity. For this, he rolls a \textbf{charm} roll and adds one grace point for each success.

\subsection{Ceremonial Service}
\index{Ceremonial Service}

Ceremonial service is a service to the deity to strengthen their work and spread their word. The service can be both a classical ceremony in memory of the deity and something like a ritual funeral or exorcism. Minor petitions or invocations may be made as part of the ceremonial service, but they do not have to be.

A ceremonial service earns the priest one grace point for each participant, up to the double \textbf{faith level} per service. If a petition or invocation is performed, this cost is deducted again.
\wbif{tirakan}{
\section{The Gods}\label{sec:die-gotter}

Once, eleven siblings awoke, who populated the world in an age long before humans, elves, and dwarves. They were daughters and sons of unknown parents, and they lived in unity until a great betrayal. Seven traitors among them broke away, leaving four gods on the side of light.

\subsection{Chronar - God of Time}\label{subsec:chronar-gott-der-zeit}

Chronar is the eldest of the siblings and watches over the flow of time. He is often depicted as the supreme god, and the largest temples and cathedrals are built for him.

Chronar is often portrayed as an old man with an hourglass, although there is no single canonical depiction. The greatest holy day associated with Chronar is the last day of the year, the Festival of the Turning of the Year.

\subsection{Nadal - Guardian of Magic}\label{subsec:nadal-wachterin-der-magie}

Nadal, the youngest of the siblings, brought magic into the world and guards it to this day. All magic in the world flows through her hands and is guided by her, or so Nadal's priests and followers preach. Nadal's temples are places of study and the teaching of magic, where peace and seclusion are valued more highly than splendor and grand services.

\subsection{Algor - Keeper of Light}\label{subsec:algor-behuter-des-lichts}

Depictions of Algor as a god are very rare; he is more often represented as the symbol of the sun in writings and on temples. He is rarely pictured as a young man holding a crystal staff, and even more rarely sketched as an old man with a simple wooden staff. These depictions date chiefly from the period between the second and sixth centuries.

Algor is celebrated on the first day of the new year with the Festival of the New Year, which immediately follows the celebrations of the Turning of the Year.

\subsection{Herbarin - The Fallen}\label{subsec:herbarin-der-gefallene}

Herbarin was the second-youngest of the eleven siblings, chosen by unknown parents to be the gods of the world of Tirakan. In the Third Age, he was mortally struck by the demon Brahas and, dying, filled with hatred and bitterness, sank into the ice of the Plains of Kaal, where all trace of him was lost forever.

\section{The Traitors}\label{sec:die-verrater}

Seven siblings broke away from the gods in the First Age and turned against them and the beings that existed in the world. Envy and resentment led them down a dark path, and from then on they pursued their own dark interests.

\subsection{Thzularn - Dark Goddess of Hatred}\label{subsec:thzularn-dunkle-gottin-des-hasses}

Thzularn is considered the instigator of the conflict between the gods. She led the dark gods and fanned the flames of hatred. She is regarded as unpredictable and unfathomably cruel, and her followers often sacrifice more than their lives to her. In return, she offers them great power over people, and still more if they are willing to give enough. Her brother and lover is Dogan, master of black magic.

\subsection{Dogan - Dark God of Magic}\label{subsec:dogan-dunkler-gott-der-magie}

Dogan, Thzularn's brother and lover, brought forth cruel, dark magic during the Demon Marches, multiplying the demons' powers and threatening the balance of the world. Though his powers are inferior to Nadal's, his strength is still present, and his followers know how to use it.

The higher arts of magic in particular tempt people to use Dogan's powers. Yet once a Gifted person has chosen to serve Dogan, he will likely never grant them freedom again.

\subsection{Seth'Nra - Dark Goddess, Guardian of the Gate}\label{subsec:sethnra-dunkle-gottin-wachterin-der-pforte}

Seth'Nra receives the dead should they choose the wrong path in the presence of Rhodd'hrom. The dead become her slaves forever, as well as slaves to the necrologists and Gifted who have chosen her path.

Old tales of terror often speak of Uhruk, Seth'Nra's driver of souls. Mercilessly, he drives those under his command to work and is said to take disobedient children into his ranks at night \ldots{}

\subsection{Rhodd'hrom - Dark God of Cursed Souls}\label{subsec:rhodd'hrom-dunkler-gott-der-verfluchten-seelen}

Long ago, when the demons had long since been driven into flight and peace was slowly returning to the world, Rhodd'hrom managed to take his place at the scales of souls, thereby determining the fate of the dead. This is the trial every person must face after death; if their soul cannot withstand the strength of the dark god, they become Seth'Nra's slave forever.

\subsection{Ranorh - Dark God of the Damned Host}\label{subsec:ranorh-dunkler-gott-des-verdammten-heeres}

When the demons were driven into the Underworld, Ranorh's influence on the world also diminished. The god, who as leader of the demons is himself wont to take the form of a demon, has a considerable following among demonologists. To be Ranorh's servant promises loyal and powerful slaves, but also a certain eternity in the Underworld.

\subsection{Ishaa'Nra - Dark Goddess of Seduction}\label{subsec:ishaa'nra-dunkle-gottin-der-verfuhrung}

Ishaa'Nra is not as powerful as her brother Chronar, and her attempts to permanently disrupt the fabric of time seem so far to have been in vain. Yet she offers her human worshippers the power to influence time, if only to a small degree.

Her arts of seduction are unsurpassed, and she also grants her followers a special charisma. Nughlhuu often appears as Ishaa'Nra's messenger; he is considered the timeless herald, the proclaimer of unnatural, premature death.

\subsection{Dhas'Garyll - The Fugitive}\label{subsec:dhasgaryll-der-geflohene}

Dhas'Garyll's whereabouts are an unsolved mystery that occupies scholars to this day. The dark god, once at the center of demonic influence, withdrew from the known planes before the demons' final banishment. His current location or condition is unknown, but he is believed to continue to exist in the world. The form he inhabits appears to be a mortal body, allowing him to act unnoticed among mortals.

\section{The Titans}\label{sec:die-titanen}

After Ranorh created thirteen cruel demons to populate the world, the gods created the Titans to protect the world of Tirakan from them. The Titans are tangible to Tirakan's people and creatures; they walk the world. They have often been observed in their deeds.

\subsection{Ce-Nya - Titaness of Mist}\label{subsec:ce-nya-titanin-des-nebels}

Ce-Nya is perhaps the most mysterious of the Titanesses. There are many rumors about her, and much remains entirely unknown. No form is associated with her; only her servants, the wolves, are encountered very often. Ce-Nya holds dominion over death, sleep, and forgetfulness, and is the patron Titaness of thieves and rogues. Her temples are hidden and often difficult for strangers to a city to recognize.

\subsection{Dilae - Titaness of Magic}\label{subsec:dilae-titanin-der-magie}

After Ce-Nya, Dilae is the most mysterious of the Titans. She serves Nadal in matters of magic, and her temples, like Nadal's, resemble libraries more than ceremonial places. Her servant creatures are perhaps the strangest beings in the world. Bolde come in every conceivable form and are unpredictable by nature. Dilae holds dominion over magic and is the patron Titaness of the Gifted.

\subsection{Ginae - Titaness of Water}\label{subsec:ginae-titanin-des-wassers}

Ginae, Titaness of Water, was created by Chronar and imprisoned the Morgalas forever in the depths of the Schlund near the Plain of Jaar. Ginae is considered exceedingly capricious and changeable in her actions in the world. She is the patron Titaness of sailors and holds dominion over birth, new beginnings, and renewal.

Her servant creatures are the hydras, those many-headed beings that dwell in the depths of all Tirakan's seas. Foremost among them is Hydonia, the legendary great hydra from the Felsenmeer.

\subsection{Gronar - Titan of Earth}\label{subsec:gronar-titan-der-erde}

Little is known of Gronar as well. Gronar serves humus and earth, and is regarded as a gentle Titan devoted to humankind. His servants are the giants, and though people very rarely catch sight of them, there are likely more than one might assume. Gronar holds dominion over life itself and is the patron Titan of farmers.

\subsection{Jogran - Titaness of Ice}\label{subsec:jogran-titanin-des-eises}

Jogran stands for ice and winter. It is said that Jogran has an ice palace far in the north, in a land inhabited by cruel, great beings. It is also said to be the origin of her servants, the unicorns. She holds dominion over winter and ice and is the patron Titaness of hunters.

\subsection{Nilan - Titan of Light}\label{subsec:nilan-titan-des-lichts}

Nilan is the Titan of Light, and people build him great, bright temples similar to those of Algor. He is served by winged serpents, native to the far reaches of the sea and only rarely seen over land. Nilan is associated with justice and law and is the patron Titan of all paladins.

\subsection{Pailos - Titan of Air}\label{subsec:pailos-titan-der-luft}

Pailos, lord of air, weather, and winds, is considered a benevolent being who, in service to the gods, guides and stirs the world's winds. He holds dominion over weather and changes in weather and is the patron Titan of jesters, con artists, and charlatans. His servants are the pegasi, winged horses found on the great northern steppes.

\subsection{Rogal - Titan of Ore}\label{subsec:rogal-titan-des-erzes}

Rogal is the Titan of Ore and the patron Titan of smiths. The Minotaurs serve his son Taurus, humanoid beings with heads like those of bulls. The Minotaurs live deep beneath the surface, and are said to be the only beings in contact with the cursed Morgalas. Rogal is associated with steadfastness and loyalty.

\subsection{Tador - Titan of Stone}\label{subsec:tador-titan-des-steins}

The mighty Tador, patron Titan of miners, is the father of stone. He holds dominion over safety and protection. Tador's servant creatures are the mighty stone beings Tonarr and Astaron, who fell in the year 722 EC in battle against Telatoon during the true Storm of Hooves.

The earthly centaurs also serve him, great beings with the bodies of horses and the torsos and heads of humans. Like the Minotaurs, centaurs live mostly underground. They rarely venture to the surface. The entrances to their caves are hidden in the forests of the north.

\subsection{Zyral - Titaness of Fire}\label{subsec:zyral-titanin-des-feuers}

Zyral brought fire to the beings of the world and armed them for battle. Created by Chronar, Zyral is the patron Titaness of warriors and fighting classes, and she holds dominion over weapons and success in battle.

Griffins serve Zyral at her command: great feathered beings with the bodies of lions and the heads of eagles. They follow Fessaal, the fire griffin from the Kemal Mountains. His mate is Phidia, the fire griffin. The Order of Phidia was founded in her honor, with the motto ``Fire, Flame, Sword''.

\section{Archdemons}\label{sec:erzdamonen2}

Created by Ranorh in the Second Age, the demons claimed the world for themselves for a long time. Only through Dhas'Garyll's actions were they driven into the demonic plane in the Third Age. These Titan-like beings now populate the Underworld.

Archdemons have direct servant creatures, usually found in their vicinity but also sent by them into the world.

\subsection{Duglaraan - The Cursed Drowner - I}\label{subsec:duglaraan-der-verfluchte-ersaufer}

The first domain is the world of water, over which Duglaraan rules with utmost cruelty. Duglaraan, the cursed drowner and defiler of seafaring, takes the form of an immense sea serpent with seven heads and great horns along the backs of its vertebrae. His two servant creatures, Shuugai and Ganod, are no better than he is. They are his companions and have a very similar form, though they are much smaller.

\subsection{Calysporn - The Cursed Ruler of Mist - II}\label{subsec:calysporn-der-verfluchte-herrscher-des-nebels}

Calysporn, the cursed ruler of mist and night, reigns over this domain. Calysporn has no fixed form and is accompanied by Dugai and Sarag'Roh, both of whom have humanoid forms, though their bodies consist only of mist.

Dugai appears male, while Sarag'Roh is likely his companion. They can slip beneath doors and through windows at night, and popular tales suspect them of killing children in their cradles. They can grant demonologists the gift of invisibility or a body of mist.

\subsection{Gyral - The Cursed Igniter - III}\label{subsec:gyral-der-verfluchte-entzunder}

The third domain is the world of fire. Gyral, also called Gyralph, rules this world with his terrible servants. Gyralph, a great figure like a troll, bears the head of an eagle. In place of a beak are three long trunks reaching to the ground. Gyralph always seems to burn; his skin resembles glowing black ash covered in countless fiery wounds. His servants, Nugshu and Tulhu, resemble griffins, but bear four horns on their backs and the same burning skin as their master.

\subsection{Cthupheen - Death from the Earth - IV}\label{subsec:cthupheen-der-tod-aus-der-erde}

Whoever encounters Cthupheen will never forget it for the rest of their life. The gigantic earthen worm measures many hundreds of paces, and its mouth bristles with thousands of shark-like teeth. Although his companions, Ijothu and Shoggu Hai, are much smaller, they still measure several paces in length and well over half a pace in diameter. The black worms are regarded as death from the earth and devour or strangle their victims.

\subsection{Brahas - Father of Addiction - V}\label{subsec:brahas-vater-der-sucht}

Brahas rules the fifth domain; he stoned the dying god Herbarin. He is called the father of addiction and the defiler of pleasures. Brahas's servants are Yoghas and Bhulu, two petrifying beings in the shape of trolls with dead, black eyes. In his vicinity, all living beings grow sluggish, lazy, and sleepy. Hardly any being, save priests and consecrated or holy animals, can escape it. Brahas's breath settles around those near him like the smoke of intoxicating poppies.

Despite their stony form and apparent lack of life, Yoghas and Bhulu are very powerful entities. They can place anyone they touch into an eternal, almost sleep-like intoxication, in which the victim is tormented by nightmares until death comes after several aeons. Yet no human has been able to document this, so the torment may last forever.

\subsection{Lhugghu - The Cursed Lord of Catastrophes - VI}\label{subsec:lhugghu-der-verfluchte-herr-der-katastrophen}

Lhugghu, who leads this domain, is the demonic lord of wind and storm, the creator of chaos and envy. Lhugghu takes the form of a faceless black cloud, and his servants, Zzibra and Cayogbaar, accompany him as malformed ravens without wings and with black skin.

\subsection{Haiphen - The Cursed Bringer of Death by Steel - VII}\label{subsec:haiphen-der-verfluchte-bringer-des-todes-mit-dem-stahl}

The seventh domain is ruled by Haiphen, who has the human form of a warrior in black armor, but without a face. Haiphen, called the cursed lord of ore and bringer of death by steel, bears a blue-glowing sword and a burning mace with which he strikes mortal wounds. His companions are Ghyogh and Hasoo, the eyeless black bloodhounds. Many mercenaries, even those serving the Lizard Army or the Minotaurs, regard him as their patron.

\subsection{Lothcam - The Cursed Master of Shadows - VIII}\label{subsec:lothcam-der-verfluchte-beherrscher-der-schatten}

Lothcam, master of shadows and shadow creatures, takes the form of a giant amoeba or a giant worm. His servants are Thurgroth and Brasho. Lothcam follows his own plan. He desires nothing less than dominion over all Tirakan.

\subsection{Nughlhuu - Ishaa'Nra's Messenger - IX}\label{subsec:nughlhuu-der-bote-ishaa'nras}

Nughlhuu is a quadruped resembling a unicorn. Black fluid constantly drips from his gaunt, almost starved body. It is said that anyone who looks into the light emanating from the horn on his head is turned to stone forever and must suffer demonic torment while Nycalhu and Nugaa slowly devour them. Nughlhuu's two companions take the form of skeletal birds, their rotting feathers hanging from pale bones.

\subsection{Iiinyca - The Defiler of Magical Power - X}\label{subsec:iiinyca-der-schander-der-magischen-kraft}

Iiinyca, the defiler of magical power, rules the tenth domain, the demonic world of magic. Iiinyca does not appear in person and can only be sensed as incredibly powerful demonic magic that disrupts all magic around it and is perceptible even to the ungifted. Iiinyca's companions, Tur Oggu and Yogzzi, are dark, cruel Bolde who make formidable opponents for any righteous person with their sharp teeth and unholy magic.

\subsection{Nycalha - The Cursed Mistress of Perverse Desire - XI}\label{subsec:nycalha-die-verfluchte-herrin-der-perversen-lust}

The eleventh domain is led by Nycalha, who brings demonic jealousy and perverse desire. Nycalha takes the form of a tall woman with pale, sickly skin and no hair on her head. Her companions, Bradhi and Flamer, appear as a pair of twins, both beautiful women: one with flaming red hair and pale skin, the other with full, sensual black hair and southern skin.

\subsection{Nyphen - Lord of Cold Death - XII}\label{subsec:nyphen-der-herr-des-kalten-todes}

Nyphen, the cursed lord of ice and cold death, is a crystalline, deep-blue figure with thousands of thorns and spikes. The gigantic, ice-cold entity is accompanied by Lhuu and Baarhai, black, toad-shaped demons. Their spikes, shot from their bodies, tear terrible wounds that never heal.

\subsection{Wisgu - Master of the First Covenant - XIII}\label{subsec:wisgu-gebieter-des-ersten-bundes}

Wisgu, hater of elves and dwarves, master of the First Covenant, is an intangible figure of dust and pure energy. In the Third Age, Wisgu made a covenant with the Morgalas and enslaved them to everlasting service. His servants Turzzi and Ylvath take the form of earthen statues.

\section{The Wayfarers}\label{sec:die-reisenden}

The gods had not yet become aware of themselves when four of these primeval beings appeared in the world. No one knew where they came from. Only they themselves could determine the purpose of their existence. The gods noticed them only late, after they had observed the world for some time.

They were four Wayfarers who came from afar. Their power equaled that of the gods, and they would decisively shape Tirakan's history in the ages to come.

\subsection{Hagarun - She Who Bears the Love and Hatred of an Entire People}\label{subsec:hagarun-die-die-liebe-und-den-hass-eines-ganzen-volkes-in-sich-tragt}

Hagarun is the Wayfarer of emotions, a being who absorbs the deepest feelings of mortals. Her appearance is mutable and reflects the souls of those who behold her. At times she appears as a radiant figure, at others as a shadowy one, depending on the emotions she carries within herself.

Her influence is subtle yet ever-present. She forges bonds between people, strengthens alliances, and kindles enmities. Many of Tirakan's wars and great love stories are attributed to her, as though she had guided the hearts of those involved. Her power lies not in direct manipulation, but in amplifying what already slumbers within mortals.

\subsection{Irminar - He Who Is Stillness}\label{subsec:irminar-der-ruhe-ist}

Irminar embodies silence and immutability as a counterpoint to Modorde's constant movement. He is the Wayfarer of balance, whose presence is felt as an endless night. In the few accounts that describe him, he is portrayed as a silent being shrouded in mist, whose eyes seem to penetrate eternity.

His power lies in stabilizing and preserving. While Modorde initiates change, Irminar ensures that the foundation remains intact. His work is most evident in moments when chaos threatens to tear the world apart. Many see him as the invisible hand that keeps Tirakan in balance.

\subsection{Modorde - He Who Ends and Recreates Being}\label{subsec:modorde-der-das-sein-beendet-und-neu-erschafft}

Modorde is a Wayfarer of change, destruction, and rebirth. His nature is marked by a deep understanding of passing away and new beginnings. He often appears as a shadowy figure whose form changes constantly, a reflection of his nature. His presence is accompanied by a sense of ending that awakens fear and hope alike.

He is not a cruel destroyer, but an inexorable renewer. To Modorde, every end is the beginning of a new cycle. His influence is especially felt in times of great upheaval, when old structures collapse and make room for the new. Legend tells that Modorde himself ended the Age of Dragons to make room for the peoples.

\subsection{Odovacar - He Who Brings Fire and Life from Afar}\label{subsec:odovacar-der-feuer-und-leben-aus-der-ferne-bringt}

Odovacar is the Wayfarer of fire, an untamed spirit who brings both life and destruction. He is often described as a burning figure, a walking storm of flames and heat. Where he walks, both fertile soils and devastated regions arise, a symbol of his dual nature.

\section{The Dragons}\label{sec:die-drachen}

Born as the children of the Wayfarers, the dragons are perhaps Tirakan's most powerful independent beings. The Wayfarers came into the world at a time when the gods were not yet even aware of their own origin, and watched events without intervening in the struggles of primordial times. At the end of the First Age, the Wayfarers begot seven children to watch over the battles between Titans and demons. These children were henceforth called dragons. According to legend, four of these seven children still live on Tirakan today, although only a few unconfirmed accounts tell of them.

\subsection{Tar - The Brave}\label{subsec:tar-der-mutige}

Tar the Brave was the first child of the Wayfarers and thus the eldest of the seven siblings. Tar, who is also the mightiest of the dragons, still dwells mainly in the northern mountains. Worshipped as a demigod by the elven peoples, he occupies the place of the highest god among the Atiarel. A river flowing from the Atiarel Mountains through the Plains of Meridian to the Felsenmeer is also named after the dragon Tar.

\subsection{Aspersia - The Fallen}\label{subsec:aspersia-die-gefallene}

During the battle at the Schlund, Aspersia was pulled down with the demons. Her fate is unknown. According to legend, however, Aspersia attempted to create a people of her own. Given the tragic circumstances and demonic surroundings, a terrible mistake must have occurred. Thus Aspersia created two peoples in her own image: one in humanoid form, the other resembling her own form. They were henceforth called the Sethlarn.

\subsection{Derumir - The Red}\label{subsec:derumir---der-rote}

Derumir, ``the Red,'' is one of the seven primordial dragons created by the Wayfarers to watch over the world. His radiant red scales and the fiery aura surrounding him make him a symbol of elemental primal power.

He appears only in times of dire need, when the very foundations of Tirakan are threatened, and his appearance is seldom a blessing for mortals. Wherever he walks, fire follows as a sign of his inseparable bond with the world's volcanic and chaotic forces. His influence is usually indirect, but perceptible in the regions where he appears. He favors the vicinity of bubbling volcanoes and primordial lava flows, and his presence can transform entire regions.

\subsection{Saranor - The Swift}\label{subsec:saranor-der-schnelle}

Saranor the Swift withdrew from Tirakan's surface following a dispute with the Wayfarers.

\subsection{Tedasiél}\label{subsec:tedasiel}

Tesasiél quarreled with Irminar during the Age of the Servant Creatures and was killed by him. Legends tell that his death was connected to the dragons' rebellion against the Wayfarers.

\subsection{Samusa - The Mother}\label{subsec:samusa-die-mutter}

Hagarun's eldest daughter is the mother of the elven people. When the dragons betrayed their parents, she was the one who gave birth to the elves. Since then, Samusa's Hoard has been spoken of in connection with the rite, among other places in the Iana Alyaria.

\subsection{Héroda - He Who Comes from the Night}\label{subsec:heroda-der-aus-der-nacht-kommt}

Héroda is one of Tirakan's seven primordial dragons, created by the Wayfarers at the end of the First Age. Like his siblings, Héroda is a divine, immortal being who keeps apart from the world of mortals. He appears only in rare moments when the balance of creation itself is threatened. Unlike Tar the Brave or Aspersia, whose fall into the Schlund is legendary, Héroda avoids every form of direct intervention. His existence is a hidden constant, a silent watch over the shadowed side of the world.

\section{Independent Entities}\label{sec:unabhangige-entitaten}

There are many independent entities of divine nature on Tirakan. Peoples such as the Quitaron or Xordai worship their own deities, unfamiliar to the realms, while entities such as Robärthrunius have a godlike status themselves, albeit without followers.

\subsection{Argorin - The Creator}\label{subsec:argorin-der-erschaffer}

The Xordai differ in their faith from most peoples of Tirakan. The latter worship, in various forms, the four or seven gods whom humans likewise regard as the highest beings in the world. The Xordai, however, pray to the two gods Argorin and Dargumir.

Argorin is the reverent one, who respects life in all its forms. He is the Creator, who keeps the embers in the heart of the world alive. His deeds are wise, and he knows the path to a better world without the terrible Sethlarn.

\subsection{Dargumir - The Fool}\label{subsec:dargumir-der-narr}

Dargumir is the Fool. He is Argorin's companion and accompanies him on the eternal quest. Unlike Argorin, he is far from inner peace. His actions are chaotic and his works sometimes cruel. He is hopelessness, but also represents life in the absolute present. Destruction and death are attributed to Dargumir.

\subsection{Akhrosch - God of the Orcs}\label{subsec:akhrosch-der-gott-der-orks}

The first father was Akhrosch. Back then, shortly after the Allfather and Allmother separated light from darkness, he appeared on Tirakan. He was called the Unborn, for he had neither father nor mother. From himself and the clay of the earth, he created Kchtuh, his companion. Together they hunted the living beings of Tirakan in those days and bore many children.

\ldots{} by Edis Elbenfreund, anthropologist of Yavon, c. 300 EC

\subsection{Taxaros - Lord of Fire}\label{subsec:taxaros-herr-des-feuers}

Taxaros is a powerful fire elemental who was brought to Tirakan, to the Academy of Jaar, by Master Astrum Fabula, probably in the eighth century. It is unclear, however, how Fabula could have accomplished this, as he was a master of black magic in Jaar. It is assumed that he had help from an elementalist in the ritual, which lasted several decades and was necessary to bind Taxaros.

\subsection{Robärthrunius - Mystical Klabauter}\label{subsec:robarthrunius-mystischer-klabauter}

Robärthrunius is a mythical king of the Klabauter who is said to sail the seas aboard an enchanted ship. Many call him godlike, but stories repeatedly suggest his mortality. Some accounts tell of wounds that even he could not heal immediately, or moments of weakness that made him pause briefly.

No one knows exactly why Robärthrunius roams the seas without rest or destination. Some say he is searching for something lost; others believe he is fleeing a great danger. Perhaps he shares more than a love of the sea with the storms; perhaps he is himself a part of them.

\subsection{Hirrukh - He Who Holds the Wheel of Creation in His Hands}\label{subsec:hirrukh-der-der-das-rad-der-schopfung-in-handen-halt}

The goblins of the Asch-Ta-Khi worship the god-being Hirrukh, who in their eyes holds the Wheel of Creation in his hands. Hirrukh is worshipped differently by the various clans: sometimes as a healer and giver of life, sometimes as a wild warrior or vengeful destroyer.

\subsection{Caltae - Minotaur Shamaness}\label{subsec:caltae-minotaurenschamanin}

Caltae is Taurus's companion. She was before the Minotaurs fell into lethargy for a thousand years, and she will be forever. Caltae lay with Taurus in the Chamber of Sleep, paralyzed and unable to move or feel emotions. When Taurus awoke in the year 0 EC, she awoke as well, and for the next 1,000 years she was to be Taurus's conscientious side.

\subsection{The Chronicler}\label{subsec:der-chronist}

The Chronicler is a mythical figure of Tirakan who is said to record the history of the world in its entirety. Some speak of him as an entity beyond this world; others claim that he is Chronar himself, preserving history in his eternal book. Alongside the Chronicler, there is also the philosophical concept of a Narrator.

The people of Tirakan know nothing of the Chronicler. The figures of the Chronicler and the Narrator exist, at most, in magico-theoretical or religious works by great philosophers. They play no role whatsoever in the events of the first 950 years of the Fifth Age.

\subsection{The Dreamer}\label{subsec:der-traumer}

He could scarcely hide his joy. He had waited a long time for this moment. For all those countless years, the Chronicler had believed him to be his son, but he knew better. He had not been born to be a Chronicler. No, he was a Dreamer.

While his ``father'' recorded the history of Tirakan as it happened, he looked over his shoulder and dreamed of how it might continue. How surprised he was when he discovered that, from time to time, one of his spiteful little ideas found its way from his dreams into the world of Tirakan. It happened quite rarely: a sacrilege here, a demonic pact there, trifles in his eyes. Yet it was a sign of his destiny.
}{}
\end{multicols}
